\chapter{Alto Paraguay}\label{dept:Alto Paraguay}
\mapaparaguai{17}{Alto Paraguay}
\vfill\clearpage
\section{Alto Paraguay}

\begin{itemize}
\tightlist
\item
  Fuerte Olimpo e a capital, ``Porta do Pantanal''; 3.471 hab.; maxima
  privacidade e suporte institucional.
\item
  Vocacao de pecuaria extensiva, pesca e logistica fluvial.
\item
  Lacuna de solar/\allowbreak pluviometria da capital fechada com dados NASA POWER
  (2001-2020).
\end{itemize}

\subsection{Ficha climática de Fuerte Olimpo (capital
departamental)}

\textbf{Irradiância solar global --- (kWh/\allowbreak m²/\allowbreak dia):}

{\def\LTcaptype{none} % do not increment counter
\begin{longtable}[]{@{}
  >{\raggedright\arraybackslash}p{(\linewidth - 22\tabcolsep) * \real{0.0833}}
  >{\raggedright\arraybackslash}p{(\linewidth - 22\tabcolsep) * \real{0.0833}}
  >{\raggedright\arraybackslash}p{(\linewidth - 22\tabcolsep) * \real{0.0833}}
  >{\raggedright\arraybackslash}p{(\linewidth - 22\tabcolsep) * \real{0.0833}}
  >{\raggedright\arraybackslash}p{(\linewidth - 22\tabcolsep) * \real{0.0833}}
  >{\raggedright\arraybackslash}p{(\linewidth - 22\tabcolsep) * \real{0.0833}}
  >{\raggedright\arraybackslash}p{(\linewidth - 22\tabcolsep) * \real{0.0833}}
  >{\raggedright\arraybackslash}p{(\linewidth - 22\tabcolsep) * \real{0.0833}}
  >{\raggedright\arraybackslash}p{(\linewidth - 22\tabcolsep) * \real{0.0833}}
  >{\raggedright\arraybackslash}p{(\linewidth - 22\tabcolsep) * \real{0.0833}}
  >{\raggedright\arraybackslash}p{(\linewidth - 22\tabcolsep) * \real{0.0833}}
  >{\raggedright\arraybackslash}p{(\linewidth - 22\tabcolsep) * \real{0.0833}}@{}}
\toprule\noalign{}
\begin{minipage}[b]{\linewidth}\raggedright
Jan
\end{minipage} & \begin{minipage}[b]{\linewidth}\raggedright
Fev
\end{minipage} & \begin{minipage}[b]{\linewidth}\raggedright
Mar
\end{minipage} & \begin{minipage}[b]{\linewidth}\raggedright
Abr
\end{minipage} & \begin{minipage}[b]{\linewidth}\raggedright
Mai
\end{minipage} & \begin{minipage}[b]{\linewidth}\raggedright
Jun
\end{minipage} & \begin{minipage}[b]{\linewidth}\raggedright
Jul
\end{minipage} & \begin{minipage}[b]{\linewidth}\raggedright
Ago
\end{minipage} & \begin{minipage}[b]{\linewidth}\raggedright
Set
\end{minipage} & \begin{minipage}[b]{\linewidth}\raggedright
Out
\end{minipage} & \begin{minipage}[b]{\linewidth}\raggedright
Nov
\end{minipage} & \begin{minipage}[b]{\linewidth}\raggedright
Dez
\end{minipage} \\
\midrule\noalign{}
\endhead
\bottomrule\noalign{}
\endlastfoot
6.43 & 5.94 & 5.66 & 4.87 & 3.73 & 3.33 & 3.58 & 4.35 & 4.83 & 5.51 &
6.31 & 6.42 \\
\end{longtable}
}

Média anual: 5.08 kWh/\allowbreak m²/\allowbreak dia. Inclinação recomendada: 21° Norte (anual);
31° inverno; 11° verão.

\textbf{Precipitação --- (mm/\allowbreak mês):}

{\def\LTcaptype{none} % do not increment counter
\begin{longtable}[]{@{}
  >{\raggedright\arraybackslash}p{(\linewidth - 22\tabcolsep) * \real{0.0833}}
  >{\raggedright\arraybackslash}p{(\linewidth - 22\tabcolsep) * \real{0.0833}}
  >{\raggedright\arraybackslash}p{(\linewidth - 22\tabcolsep) * \real{0.0833}}
  >{\raggedright\arraybackslash}p{(\linewidth - 22\tabcolsep) * \real{0.0833}}
  >{\raggedright\arraybackslash}p{(\linewidth - 22\tabcolsep) * \real{0.0833}}
  >{\raggedright\arraybackslash}p{(\linewidth - 22\tabcolsep) * \real{0.0833}}
  >{\raggedright\arraybackslash}p{(\linewidth - 22\tabcolsep) * \real{0.0833}}
  >{\raggedright\arraybackslash}p{(\linewidth - 22\tabcolsep) * \real{0.0833}}
  >{\raggedright\arraybackslash}p{(\linewidth - 22\tabcolsep) * \real{0.0833}}
  >{\raggedright\arraybackslash}p{(\linewidth - 22\tabcolsep) * \real{0.0833}}
  >{\raggedright\arraybackslash}p{(\linewidth - 22\tabcolsep) * \real{0.0833}}
  >{\raggedright\arraybackslash}p{(\linewidth - 22\tabcolsep) * \real{0.0833}}@{}}
\toprule\noalign{}
\begin{minipage}[b]{\linewidth}\raggedright
Jan
\end{minipage} & \begin{minipage}[b]{\linewidth}\raggedright
Fev
\end{minipage} & \begin{minipage}[b]{\linewidth}\raggedright
Mar
\end{minipage} & \begin{minipage}[b]{\linewidth}\raggedright
Abr
\end{minipage} & \begin{minipage}[b]{\linewidth}\raggedright
Mai
\end{minipage} & \begin{minipage}[b]{\linewidth}\raggedright
Jun
\end{minipage} & \begin{minipage}[b]{\linewidth}\raggedright
Jul
\end{minipage} & \begin{minipage}[b]{\linewidth}\raggedright
Ago
\end{minipage} & \begin{minipage}[b]{\linewidth}\raggedright
Set
\end{minipage} & \begin{minipage}[b]{\linewidth}\raggedright
Out
\end{minipage} & \begin{minipage}[b]{\linewidth}\raggedright
Nov
\end{minipage} & \begin{minipage}[b]{\linewidth}\raggedright
Dez
\end{minipage} \\
\midrule\noalign{}
\endhead
\bottomrule\noalign{}
\endlastfoot
122 & 134 & 100 & 78 & 80 & 40 & 26 & 16 & 41 & 99 & 141 & 128 \\
\end{longtable}
}

Média anual: \textasciitilde1.005 mm/\allowbreak ano. Regime sub-úmido Chaco Boreal.
Estação chuvosa Nov-Fev; seca Jun-Set.
\clearpage
\subsection*{Dados Departamentais}
\subsubsection{Desenvolvimento Humano e
Social}

\subsection{IDH Departamental}

{\def\LTcaptype{none} % do not increment counter
\begin{longtable}[]{@{}
  >{\raggedright\arraybackslash}p{(\linewidth - 2\tabcolsep) * \real{0.6111}}
  >{\raggedright\arraybackslash}p{(\linewidth - 2\tabcolsep) * \real{0.3889}}@{}}
\toprule\noalign{}
\begin{minipage}[b]{\linewidth}\raggedright
Indicador
\end{minipage} & \begin{minipage}[b]{\linewidth}\raggedright
Valor
\end{minipage} \\
\midrule\noalign{}
\endhead
\bottomrule\noalign{}
\endlastfoot
IDH & 0,619 \\
Classificação IDH & baixo \\
Ranking nacional (1=melhor) & 18/\allowbreak 18 (ou 17/\allowbreak 18, disputado com Caazapá) \\
Esperança de vida & 68,4 anos \\
Anos de escolaridade média & 5,8 anos \\
RNB per capita (USD PPA) & 4.900 \\
\end{longtable}
}

\subsection{Indicadores de Pobreza}

{\def\LTcaptype{none} % do not increment counter
\begin{longtable}[]{@{}ll@{}}
\toprule\noalign{}
Indicador & Valor \\
\midrule\noalign{}
\endhead
\bottomrule\noalign{}
\endlastfoot
Taxa de pobreza total (\%) & 38,7\% \\
Taxa de pobreza extrema (\%) & estimativa 14,2\% \\
Índice de Gini & estimativa 0,478 \\
\end{longtable}
}

\subsection{Infraestrutura Básica}

{\def\LTcaptype{none} % do not increment counter
\begin{longtable}[]{@{}ll@{}}
\toprule\noalign{}
Indicador & Valor \\
\midrule\noalign{}
\endhead
\bottomrule\noalign{}
\endlastfoot
Acesso a água potável (\%) & estimativa 34,5\% \\
Acesso a saneamento (\%) & estimativa 31,2\% \\
\end{longtable}
}

\subsubsection{Segurança Pública}

\subsection{Indicadores}

{\def\LTcaptype{none} % do not increment counter
\begin{longtable}[]{@{}
  >{\raggedright\arraybackslash}p{(\linewidth - 2\tabcolsep) * \real{0.6111}}
  >{\raggedright\arraybackslash}p{(\linewidth - 2\tabcolsep) * \real{0.3889}}@{}}
\toprule\noalign{}
\begin{minipage}[b]{\linewidth}\raggedright
Indicador
\end{minipage} & \begin{minipage}[b]{\linewidth}\raggedright
Valor
\end{minipage} \\
\midrule\noalign{}
\endhead
\bottomrule\noalign{}
\endlastfoot
Taxa de homicídios (por 100k hab) & \textasciitilde2--4 (estimativa,
abaixo da média nacional) \\
Crimes registrados (total/\allowbreak ano) & muito baixo (população total
\textasciitilde15 mil hab.) \\
Número de comisarías & \textasciitilde3 (estimativa --- população mínima
em território imenso) \\
Índice segurança relativo & alto \\
\end{longtable}
}

\subsubsection{Mercado de Terra Rural}

\subsection{Preços por Tipo de Terra}

{\def\LTcaptype{none} % do not increment counter
\begin{longtable}[]{@{}llll@{}}
\toprule\noalign{}
Tipo & Preço (USD/\allowbreak ha) & Faixa & Tendência \\
\midrule\noalign{}
\endhead
\bottomrule\noalign{}
\endlastfoot
Agrícola (limitada, zonas mais elevadas) & 1.200 & 800 -- 2.000 & → \\
Pastagem para pecuária & 500 & 300 -- 900 & ↑ \\
Terra bruta /\allowbreak  pantanosa & 200 & 100 -- 400 & → \\
\end{longtable}
}

\subsection{Características do
Mercado}

{\def\LTcaptype{none} % do not increment counter
\begin{longtable}[]{@{}ll@{}}
\toprule\noalign{}
Parâmetro & Valor \\
\midrule\noalign{}
\endhead
\bottomrule\noalign{}
\endlastfoot
Valorização anual média (\%) & 3--5\% \\
Lote mínimo rural (ha) & 100 (viável apenas em grandes escalas) \\
Imposto territorial rural (\%) & 1\% sobre valor fiscal \\
Liquidez do mercado & muito baixa \\
\end{longtable}
}
\clearpage
\newpage
\secaoDiagnostico{Bahia Negra}{\mapaDistrito{1704}}{Bahía Negra é o ``fim do mundo'' estratégico no Paraguai. Oferece o
nível mais radical de isolamento geográfico e demográfico do país,
transformando a dificuldade de acesso em sua maior fortaleza. Sua força
absoluta reside no anonimato total garantido pelo vazio populacional e
pela imensidão do território. Contudo, essa mesma característica impõe
desafios brutais de logística: a sobrevivência aqui exige um perfil
anfíbio, com autonomia total para viver sem suporte do Estado por meses
a fio. É o local ideal para quem busca o máximo de privacidade e
segurança através da inacessibilidade, transformando os humedales do
Pantanal em uma barreira natural contra qualquer colapso sistêmico
externo.

\subsubsection{Por que sim}

\begin{itemize}
\tightlist
\item
  Máximo isolamento geográfico e demográfico (privacidade absoluta).
\item
  Ambiente social mais seguro do país pelo distanciamento físico.
\item
  Abundância massiva de proteína animal (gado e pesca).
\item
  Distância máxima de qualquer alvo tático ou eixo de fallout tático.
\end{itemize}

\subsubsection{Por que não}

\begin{itemize}
\tightlist
\item
  Risco hidrológico extremo (cheias sazonais e isolamento prolongado).
\item
  Inexistência de asfalto e serviços de saúde especializados.
\item
  Dependência total de logística fluvial ou aérea para suprimentos
  técnicos.
\end{itemize}

\subsubsection{Recomendação
Técnica}

Indicado exclusivamente para estabelecimentos de isolamento tático
extremo que possuam infraestrutura própria de energia, saúde e
transporte (pista de pouso ou barco). Requer estocagem radical de
insumos (6 meses). O GSS de 6.5 reflete a superioridade do isolamento
protegida pela hostilidade geográfica da região. Referência atual de
score: GSS 6.5.

\subsubsection{}

\begin{itemize}
\tightlist
\item
  \begin{enumerate}
  \def\labelenumi{\arabic{enumi})}
  \tightlist
  \item
    Dados censitários validados (INE\footnote{Instituto Nacional de Estadística (Instituto Nacional de Estatística), órgão oficial de dados demográficos e censitários.} 2022).
  \end{enumerate}
\item
  \begin{enumerate}
  \def\labelenumi{\arabic{enumi})}
  \setcounter{enumi}{1}
  \tightlist
  \item
    Relatórios de monitoramento de cheias do Pantanal (SEN\footnote{Secretaría de Emergencia Nacional (Secretaria de Emergência Nacional), órgão governamental de resposta a desastres.}/\allowbreak MADES).
  \end{enumerate}
\item
  \begin{enumerate}
  \def\labelenumi{\arabic{enumi})}
  \setcounter{enumi}{2}
  \tightlist
  \item
    Mapa de solos e acessibilidade rural do Alto Paraguay (MOPC).
  \end{enumerate}
\item
  \begin{enumerate}
  \def\labelenumi{\arabic{enumi})}
  \setcounter{enumi}{3}
  \tightlist
  \item
    Registro de segurança pública e governança de fronteira naval.
  \end{enumerate}
\end{itemize}}{dist:17_Alto_Paraguay:Bahia_Negra}
\begin{table}[ht!]
\small \begin{tabular}{p{3.0cm}p{0.8cm}p{6.0cm}} \toprule
\textbf{Critério} & \textbf{Nota} & \textbf{Justificativa} \\ \midrule
A - Ameaças Estratégicas & 8.5 & Isolamento tático absoluto; ausência total de interesse estratégico militar em conflito global \\
B - Risco Social & 9.0 & A menor densidade populacional do país; segurança contra riscos de massa absoluta \\
C - Riscos Naturais & 3.5 & Risco hidrológico extremo; nota penalizada pelo isolamento cíclico por lama e cheias \\
D - Autossuficiência & 6.0 & Abundância de proteína animal e água; severamente limitado pela logística inexistente \\
E - Institucional & 4.5 & Ambiente social seguro pelo isolamento e presença militar de fronteira funcional \\
\midrule \textbf{GSS FINAL} & \textbf{6.5} & \textbf{Moderadamente Seguro (Indicado apenas para projetos de refúgio tático extremo "off-grid" que possuam autonomia logística e hídrica total).} \\ \bottomrule \end{tabular}
\end{table}
\newpage
\secaoDiagnostico{Capitan Carmelo Peralta}{\mapaConteudo{-21.6833}{-57.9000}}{Capitán Carmelo Peralta (Alto Paraguay) é um hub de fronteira em
transição radical. Deixou de ser uma zona isolada para se tornar o
epicentro da Rota Bioceânica, o que traz oportunidades massivas de
autossuficiência proteica e valorização, mas também riscos estratégicos
elevados por ser um nó logístico internacional.

\subsubsection{Por que sim}

\begin{itemize}
\tightlist
\item
  Isolamento geográfico por densidade, mas com conectividade moderna de
  alta qualidade.
\item
  Potencial pecuário e recursos hídricos significativos.
\item
  Suporte estatal prioritário por ser projeto de interesse nacional.
\end{itemize}

\subsubsection{Por que não}

\begin{itemize}
\tightlist
\item
  Alvo estratégico evidente em cenários de conflito
  logístico/\allowbreak fronteiriço.
\item
  Risco de inundação sazonal requer obras de proteção caras.
\item
  Restrições de compra de terra para estrangeiros de países limítrofes.
\end{itemize}

\subsubsection{Pre-condicoes Minimas de
Decisao}

\begin{itemize}
\tightlist
\item
  Validação jurídica de aquisição de terras (Lei 2532/\allowbreak 05).
\item
  Avaliação da cota de inundação no terreno específico.
\item
  Implementação de autonomia energética (Solar) e de água.
\end{itemize}

\subsubsection{Recomendação
Técnica}

Uso recomendado para perfis que buscam valorização estratégica e
autossuficiência em áreas de baixa densidade, cientes do papel
geopolítico do hub. Referência atual de score: GSS 6.8; risco
predominante: moderado (hidrológico e estratégico).

\subsubsection{}

\begin{itemize}
\tightlist
\item
  \begin{enumerate}
  \def\labelenumi{\arabic{enumi})}
  \tightlist
  \item
    Delimitacao territorial validada por georreferencia institucional.
  \end{enumerate}
\item
  \begin{enumerate}
  \def\labelenumi{\arabic{enumi})}
  \setcounter{enumi}{1}
  \tightlist
  \item
    População municipal referenciada por Censo 2022.
  \end{enumerate}
\item
  \begin{enumerate}
  \def\labelenumi{\arabic{enumi})}
  \setcounter{enumi}{2}
  \tightlist
  \item
    Indicadores de infraestrutura e rodovias PY15.
  \end{enumerate}
\item
  \begin{enumerate}
  \def\labelenumi{\arabic{enumi})}
  \setcounter{enumi}{3}
  \tightlist
  \item
    Riscos climáticos e hidrológicos.
  \end{enumerate}
\item
  \begin{enumerate}
  \def\labelenumi{\arabic{enumi})}
  \setcounter{enumi}{4}
  \tightlist
  \item
    Informações de custo de energia e matriz.
  \end{enumerate}
\item
  \begin{enumerate}
  \def\labelenumi{\arabic{enumi})}
  \setcounter{enumi}{5}
  \tightlist
  \item
    Restrições legais de fronteira - Lei 2532/\allowbreak 05.
  \end{enumerate}
\end{itemize}}{dist:17_Alto_Paraguay:Capitan_Carmelo_Peralta}
\begin{table}[ht!]
\small \begin{tabular}{p{3.0cm}p{0.8cm}p{6.0cm}} \toprule
\textbf{Critério} & \textbf{Nota} & \textbf{Justificativa} \\ \midrule
A - Ameaças Estratégicas & 5.0 & Hub logístico internacional; alvo estratégico em conflitos \\
B - Risco Social & 7.5 & Baixíssima densidade demográfica compensada pelo rápido crescimento \\
C - Riscos Naturais & 6.5 & Geologia estável no Chaco; vulnerabilidade hídrica sazonal \\
D - Autossuficiência & 8.0 & Elevado potencial de autossuficiência proteica e hub de suprimentos \\
E - Institucional & 7.4 & Forte suporte governamental; restrição legal de fronteira para estrangeiros \\
\midrule \textbf{GSS FINAL} & \textbf{6.8} & \textbf{Moderadamente Seguro (Exige preparação e infraestrutura dedicada).} \\ \bottomrule \end{tabular}
\end{table}
\subsubsection{Vulnerabilidades e
Mitigação}\label{vuln-Capitan_Carmelo_Peralta-vulnerabilidades-e-mitigauxe7uxe3o}

\begin{itemize}
\tightlist
\item
  \textbf{Exposicao Hidrologica:} Risco de inundação sazonal pelo Rio
  Paraguai (reforço de defesas costeiras e construção elevada
  recomendada).
\item
  \textbf{Isolamento Logístico:} Dependência de poucas rotas de saída; o
  Corredor Bioceânico é vital (monitorar tráfego e possíveis bloqueios).
\item
  \textbf{Segurança de Fronteira:} Proximidade com o Brasil exige
  vigilância contra fluxos migratórios descontrolados ou crises
  transfronteiriças.
\item
  \textbf{Autossuficiência Energética:} Baixa diversidade de fontes
  locais (potencializar microgeração solar para reduzir dependência da
  rede).
\end{itemize}

\subsubsection{Dossiê de Campo}\label{dossie-Capitan_Carmelo_Peralta-dossiuxea-de-campo}
\begin{description}[style=nextline,leftmargin=0.5cm,font=\bfseries]
\item[Coordenadas]  21°41'S 57°54'W.
\item[Topografia]  Relevo plano (geologicamente estável, baixo risco sísmico), margeado pelo Rio Paraguai (fronteira com Brasil). Área de \textasciitilde 4.798 km².
\item[Alvos Estratégicos]  Ponte Bioceânica (em fase final de construção - conexão vital Paraguai-Brasil); Eixo da Ruta PY15 (Corredor Bioceânico); Porto fluvial; Proximidade com o Pantanal paraguaio. Localização estratégica absoluta como hub logístico internacional de carga (celulose, grãos e proteína).
\item[Fallout]  Ventos predominantes do Norte/\allowbreak Nordeste.
\item[População]  4.402 (Censo 2022 Final). População em rápido crescimento devido ao impacto das obras de infraestrutura. Presença de comunidades indígenas Ayoreo.
\item[Densidade]  \textasciitilde 0,92 hab/\allowbreak km² (Extremamente baixa, oferecendo isolamento geográfico massivo fora do núcleo urbano).
\item[Vias de Acesso]  Transformação radical com o Corredor Bioceânico pavimentado (PY15) e a nova ponte internacional para Porto Murtinho (Brasil). Risco elevado de saturação e controle militar em crises devido ao valor logístico da rota.
\item[Serviços]  Infraestrutura urbana e de serviços em fase de expansão acelerada. Dependência residual de centros maiores para saúde de alta complexidade.
\item[Custo de Vida]  Em ascensão devido à valorização imobiliária; energia elétrica barata (ANDE\footnote{Administración Nacional de Electricidade (Administração Nacional de Eletricidade), companhia estatal de energia.} \textasciitilde US\$ 0,05/\allowbreak kWh).
\item[Preço da Terra]  US\$ 300-600/\allowbreak ha (bruta) a US\$ 1.200/\allowbreak ha (desenvolvida). Forte especulação em áreas logísticas.
\item[Hidrologia]  Risco de Inundação Sazonal. Localidade margeada pelo Rio Paraguai, vulnerável a cheias extraordinárias. Obras de defesa costeira em andamento. Risco de incêndios florestais em períodos de seca no Chaco.
\item[Clima]  Tropical Úmido/\allowbreak Seco. Verões com calor extremo. Déficit hídrico é o principal desafio agrícola.
\item[Sismicidade]  Risco baixo; região estável com eventos intraplaca raros e de baixa magnitude (<4.0 Richter).
\item[Energia]  Rede nacional (Itaipu) em fase de reforço para atender ao novo hub logístico.
\end{description}
\textbf{Fonte climática:} NASA POWER Climatology API (período 2001-2020)

\textbf{Fonte luminosa:} estimativa world atlas

\paragraph{Irradiação Solar
(kWh/\allowbreak m²/\allowbreak dia)}\label{irradiauxe7uxe3o-solar-kwhmuxb2dia}

{\def\LTcaptype{none} % do not increment counter
\begin{longtable}[]{@{}
  >{\raggedright\arraybackslash}p{(\linewidth - 24\tabcolsep) * \real{0.0746}}
  >{\raggedright\arraybackslash}p{(\linewidth - 24\tabcolsep) * \real{0.0746}}
  >{\raggedright\arraybackslash}p{(\linewidth - 24\tabcolsep) * \real{0.0746}}
  >{\raggedright\arraybackslash}p{(\linewidth - 24\tabcolsep) * \real{0.0746}}
  >{\raggedright\arraybackslash}p{(\linewidth - 24\tabcolsep) * \real{0.0746}}
  >{\raggedright\arraybackslash}p{(\linewidth - 24\tabcolsep) * \real{0.0746}}
  >{\raggedright\arraybackslash}p{(\linewidth - 24\tabcolsep) * \real{0.0746}}
  >{\raggedright\arraybackslash}p{(\linewidth - 24\tabcolsep) * \real{0.0746}}
  >{\raggedright\arraybackslash}p{(\linewidth - 24\tabcolsep) * \real{0.0746}}
  >{\raggedright\arraybackslash}p{(\linewidth - 24\tabcolsep) * \real{0.0746}}
  >{\raggedright\arraybackslash}p{(\linewidth - 24\tabcolsep) * \real{0.0746}}
  >{\raggedright\arraybackslash}p{(\linewidth - 24\tabcolsep) * \real{0.0746}}
  >{\raggedright\arraybackslash}p{(\linewidth - 24\tabcolsep) * \real{0.1045}}@{}}
\toprule\noalign{}
\begin{minipage}[b]{\linewidth}\raggedright
Jan
\end{minipage} & \begin{minipage}[b]{\linewidth}\raggedright
Fev
\end{minipage} & \begin{minipage}[b]{\linewidth}\raggedright
Mar
\end{minipage} & \begin{minipage}[b]{\linewidth}\raggedright
Abr
\end{minipage} & \begin{minipage}[b]{\linewidth}\raggedright
Mai
\end{minipage} & \begin{minipage}[b]{\linewidth}\raggedright
Jun
\end{minipage} & \begin{minipage}[b]{\linewidth}\raggedright
Jul
\end{minipage} & \begin{minipage}[b]{\linewidth}\raggedright
Ago
\end{minipage} & \begin{minipage}[b]{\linewidth}\raggedright
Set
\end{minipage} & \begin{minipage}[b]{\linewidth}\raggedright
Out
\end{minipage} & \begin{minipage}[b]{\linewidth}\raggedright
Nov
\end{minipage} & \begin{minipage}[b]{\linewidth}\raggedright
Dez
\end{minipage} & \begin{minipage}[b]{\linewidth}\raggedright
Média
\end{minipage} \\
\midrule\noalign{}
\endhead
\bottomrule\noalign{}
\endlastfoot
6.74 & 6.20 & 5.50 & 4.48 & 3.35 & 2.88 & 3.21 & 3.96 & 4.59 & 5.43 &
6.43 & 6.81 & \textbf{4.96} \\
\end{longtable}
}

\textbf{Inclinação solar recomendada:} 25° N (anual) · 35° N (inverno
jun-ago) · 15° N (verão nov-jan)

\paragraph{Precipitação (mm/\allowbreak mês)}\label{precipitauxe7uxe3o-mmmuxeas}

{\def\LTcaptype{none} % do not increment counter
\begin{longtable}[]{@{}
  >{\raggedright\arraybackslash}p{(\linewidth - 24\tabcolsep) * \real{0.0704}}
  >{\raggedright\arraybackslash}p{(\linewidth - 24\tabcolsep) * \real{0.0704}}
  >{\raggedright\arraybackslash}p{(\linewidth - 24\tabcolsep) * \real{0.0704}}
  >{\raggedright\arraybackslash}p{(\linewidth - 24\tabcolsep) * \real{0.0704}}
  >{\raggedright\arraybackslash}p{(\linewidth - 24\tabcolsep) * \real{0.0704}}
  >{\raggedright\arraybackslash}p{(\linewidth - 24\tabcolsep) * \real{0.0704}}
  >{\raggedright\arraybackslash}p{(\linewidth - 24\tabcolsep) * \real{0.0704}}
  >{\raggedright\arraybackslash}p{(\linewidth - 24\tabcolsep) * \real{0.0704}}
  >{\raggedright\arraybackslash}p{(\linewidth - 24\tabcolsep) * \real{0.0704}}
  >{\raggedright\arraybackslash}p{(\linewidth - 24\tabcolsep) * \real{0.0704}}
  >{\raggedright\arraybackslash}p{(\linewidth - 24\tabcolsep) * \real{0.0704}}
  >{\raggedright\arraybackslash}p{(\linewidth - 24\tabcolsep) * \real{0.0704}}
  >{\raggedright\arraybackslash}p{(\linewidth - 24\tabcolsep) * \real{0.1549}}@{}}
\toprule\noalign{}
\begin{minipage}[b]{\linewidth}\raggedright
Jan
\end{minipage} & \begin{minipage}[b]{\linewidth}\raggedright
Fev
\end{minipage} & \begin{minipage}[b]{\linewidth}\raggedright
Mar
\end{minipage} & \begin{minipage}[b]{\linewidth}\raggedright
Abr
\end{minipage} & \begin{minipage}[b]{\linewidth}\raggedright
Mai
\end{minipage} & \begin{minipage}[b]{\linewidth}\raggedright
Jun
\end{minipage} & \begin{minipage}[b]{\linewidth}\raggedright
Jul
\end{minipage} & \begin{minipage}[b]{\linewidth}\raggedright
Ago
\end{minipage} & \begin{minipage}[b]{\linewidth}\raggedright
Set
\end{minipage} & \begin{minipage}[b]{\linewidth}\raggedright
Out
\end{minipage} & \begin{minipage}[b]{\linewidth}\raggedright
Nov
\end{minipage} & \begin{minipage}[b]{\linewidth}\raggedright
Dez
\end{minipage} & \begin{minipage}[b]{\linewidth}\raggedright
Total/\allowbreak ano
\end{minipage} \\
\midrule\noalign{}
\endhead
\bottomrule\noalign{}
\endlastfoot
131 & 136 & 131 & 152 & 172 & 90 & 75 & 52 & 92 & 187 & 190 & 166 &
\textbf{1577 mm} \\
\end{longtable}
}

\paragraph{Poluição Luminosa}\label{poluiuxe7uxe3o-luminosa}

{\def\LTcaptype{none} % do not increment counter
\begin{longtable}[]{@{}ll@{}}
\toprule\noalign{}
Parâmetro & Valor \\
\midrule\noalign{}
\endhead
\bottomrule\noalign{}
\endlastfoot
Escala Bortle & 1 --- Céu verdadeiramente escuro \\
Radiância artificial & 0.1 nW/\allowbreak cm²/\allowbreak sr \\
\end{longtable}
}
\paragraph{Solo (SoilGrids 2.0, média ponderada 0--30
cm)}

{\def\LTcaptype{none} % do not increment counter
\begin{longtable}[]{@{}ll@{}}
\toprule\noalign{}
Parâmetro & Valor \\
\midrule\noalign{}
\endhead
\bottomrule\noalign{}
\endlastfoot
pH (H$_2$O) & 6.25 \\
Carbono orgânico (SOC) & 14.97 g/\allowbreak kg \\
Argila & 32.23 \% \\
Areia & 31.27 \% \\
Silte (calc.) & 36.5 \% \\
Densidade aparente & 1.31 g/\allowbreak cm³ \\
\textbf{Aptidão agrícola} & \textbf{Alta} \\
\end{longtable}
}

\subsubsection{Indicadores Sociais}

\textbf{Fontes:} PNUD 2020, INE\footnote{Instituto Nacional de Estadística (Instituto Nacional de Estatística), órgão oficial de dados demográficos e censitários.} Censo 2022, INE\footnote{Instituto Nacional de Estadística (Instituto Nacional de Estatística), órgão oficial de dados demográficos e censitários.} EPHC 2023, Ministerio
Público 2024\\
\textbf{Nota:} Valores marcados como \emph{(dept.)} referem-se ao
departamento de Alto Paraguay; valores \emph{(dist.)} são específicos
deste distrito. Dados marcados \emph{(est.)} são estimativas.

{\def\LTcaptype{none} % do not increment counter
\begin{longtable}[]{@{}
  >{\raggedright\arraybackslash}p{(\linewidth - 4\tabcolsep) * \real{0.4231}}
  >{\raggedright\arraybackslash}p{(\linewidth - 4\tabcolsep) * \real{0.2692}}
  >{\raggedright\arraybackslash}p{(\linewidth - 4\tabcolsep) * \real{0.3077}}@{}}
\toprule\noalign{}
\begin{minipage}[b]{\linewidth}\raggedright
Indicador
\end{minipage} & \begin{minipage}[b]{\linewidth}\raggedright
Valor
\end{minipage} & \begin{minipage}[b]{\linewidth}\raggedright
Âmbito
\end{minipage} \\
\midrule\noalign{}
\endhead
\bottomrule\noalign{}
\endlastfoot
IDH (2020) & 0,619 & dept. \\
Ranking IDH nacional & 18/\allowbreak 18 (ou 17/\allowbreak 18, disputado com Caazapá) &
dept. \\
Esperança de vida & 68,4 anos & dept. \\
Escolaridade média & 5,8 anos & dept. \\
RNB per capita (USD PPA) & 4.900 & dept. \\
Pobreza monetária (\%) & 38,7\% & dept. \\
Pobreza extrema (\%) & estimativa 14,2\% & dept. \\
Índice de Gini & estimativa 0,478 & dept. \\
Acesso a água potável (\%) & estimativa 34,5\% & dept. \\
Acesso a saneamento (\%) & estimativa 31,2\% & dept. \\
Taxa de homicídios (est., /\allowbreak 100k hab) & \textasciitilde2--4 (estimativa,
abaixo da média nacional) & dept. \\
Índice de segurança & alto & dept. \\
População (Censo 2022) & 4.402 & dist. \\
Idade mediana & N/\allowbreak D & dist. \\
Taxa de fecundidade & N/\allowbreak D & dist. \\
\end{longtable}
}
\subsubsection{Combustível}

\textbf{Referência:} PETROPAR /\allowbreak  postos locais (2024) \textbf{Tipo de
localidade:} Interior

{\def\LTcaptype{none} % do not increment counter
\begin{longtable}[]{@{}lll@{}}
\toprule\noalign{}
Combustível & USD/\allowbreak litro & Gs/\allowbreak litro (aprox.) \\
\midrule\noalign{}
\endhead
\bottomrule\noalign{}
\endlastfoot
Gasolina 93 oct & 1.18 & 8,732 \\
Gasolina 97 oct (premium) & 1.28 & 9,472 \\
Diesel & 1.10 & 8,140 \\
\end{longtable}
}

\begin{quote}
Preços podem variar ±5\% conforme posto e sazonalidade. Chaco e interior
remoto apresentam maior variação.
\end{quote}

\subsubsection{Cobertura Celular}

\textbf{Fonte:} CONATEL PY /\allowbreak  operadoras (2024)

{\def\LTcaptype{none} % do not increment counter
\begin{longtable}[]{@{}ll@{}}
\toprule\noalign{}
Parâmetro & Valor \\
\midrule\noalign{}
\endhead
\bottomrule\noalign{}
\endlastfoot
Cobertura 4G população (dept.) & 35\% \\
Cobertura 4G área rural & 15\% \\
Melhor operadora & Tigo \\
Qualidade rural & limitada \\
\end{longtable}
}

\begin{quote}
Para áreas rurais fora do núcleo urbano, recomenda-se chip Tigo como
principal e Personal como backup.
\end{quote}

\subsubsection{Internet}

\textbf{Fonte:} CONATEL /\allowbreak  Speedtest Ookla (2024)

{\def\LTcaptype{none} % do not increment counter
\begin{longtable}[]{@{}ll@{}}
\toprule\noalign{}
Parâmetro & Valor \\
\midrule\noalign{}
\endhead
\bottomrule\noalign{}
\endlastfoot
Velocidade média download & 15 Mbps \\
Domicílios com internet (dept.) & 22\% \\
Tecnologia predominante & satélite \\
Opção rural & Starlink disponível (\textasciitilde USD 44/\allowbreak mês) \\
\end{longtable}
}

\subsubsection{Mercado Imobiliário e Terra
Rural}

\textbf{Fonte:} INDERT /\allowbreak  Clasificados.com.py (2024)

{\def\LTcaptype{none} % do not increment counter
\begin{longtable}[]{@{}ll@{}}
\toprule\noalign{}
Tipo & Referência \\
\midrule\noalign{}
\endhead
\bottomrule\noalign{}
\endlastfoot
Terra agrícola alta prod. (USD/\allowbreak ha) & 600 \\
Imóvel urbano (USD/\allowbreak m²) & 400 \\
Aluguel 2 quartos (USD/\allowbreak mês) & 160 \\
\end{longtable}
}

\begin{quote}
Valores de referência departamental. Localidades menores podem ter
preços 20--40\% abaixo da capital departamental.
\end{quote}

\subsubsection{Saúde}

\textbf{Fonte:} MSPBS /\allowbreak  IPS Paraguay (2024-2026), consolidação
departamental e proxy local

{\def\LTcaptype{none} % do not increment counter
\begin{longtable}[]{@{}ll@{}}
\toprule\noalign{}
Serviço & Disponibilidade \\
\midrule\noalign{}
\endhead
\bottomrule\noalign{}
\endlastfoot
USF /\allowbreak  Posto de Saúde & sim \\
Hospital Regional & não \\
IPS (seguro social) & não \\
Farmácia & sim \\
Distância ao hospital de referência & \textasciitilde641 km (Fuerte
Olimpo) \\
\end{longtable}
}

\textbf{Principais estabelecimentos:} USF local; referencia hospitalar
em Fuerte Olimpo

\textbf{Observação para imigrantes:} Atendimento primario local ou em
raio curto; casos de maior complexidade seguem para Fuerte Olimpo.
Cobertura privada continua recomendada para especialidades e urgências
de maior complexidade.
\newpage
\secaoDiagnostico{Fuerte Olimpo}{\mapaDistrito{1701}}{Fuerte Olimpo é a ``fortaleza do norte'' paraguaio. Oferece o nível mais
elevado de isolamento tático e demográfico do país, protegida pela
geografia remota e por colinas que funcionam como observatórios
naturais. Sua força absoluta reside na combinação de suporte
institucional (Capital) com um vazio populacional que garante
privacidade total. O desafio principal é a logística terrestre precária,
que exige do ocupante autonomia de transporte fluvial e estoque massivo
de suprimentos. É o local ideal para quem busca segurança através da
inacessibilidade, transformando a capital de Alto Paraguay no refúgio
mais resiliente e discreto da Região Ocidental.

\subsubsection{Por que sim}

\begin{itemize}
\tightlist
\item
  Máximo isolamento geográfico e demográfico (anonimato absoluto).
\item
  Abundância massiva de proteína animal e água fluvial perene.
\item
  Ambiente social extremamente seguro sob proteção capital.
\item
  Longe de qualquer alvo militar terrestre ou eixo de fallout tático.
\end{itemize}

\subsubsection{Por que não}

\begin{itemize}
\tightlist
\item
  Risco de isolamento físico total por meses em eventos climáticos.
\item
  Inexistência de asfalto e dependência de logística fluvial/\allowbreak aérea.
\item
  Restrições da Lei de Fronteira para investidores estrangeiros.
\end{itemize}

\subsubsection{Recomendação
Técnica}

Altamente indicado para residências de isolamento tático extremo ou
bases de inteligência regional. Requer autonomia energética, hídrica e
logística total. O GSS de 8.1 reflete a excepcional força das
instituições locais e o isolamento protegida pela hostilidade geográfica
da região. Referência atual de score: GSS 8.1.

\subsubsection{}

\begin{itemize}
\tightlist
\item
  \begin{enumerate}
  \def\labelenumi{\arabic{enumi})}
  \tightlist
  \item
    Dados censitários validados (INE\footnote{Instituto Nacional de Estadística (Instituto Nacional de Estatística), órgão oficial de dados demográficos e censitários.} 2022).
  \end{enumerate}
\item
  \begin{enumerate}
  \def\labelenumi{\arabic{enumi})}
  \setcounter{enumi}{1}
  \tightlist
  \item
    Relatórios de monitoramento ambiental do Rio Paraguai (SEN\footnote{Secretaría de Emergencia Nacional (Secretaria de Emergência Nacional), órgão governamental de resposta a desastres.}/\allowbreak DMH).
  \end{enumerate}
\item
  \begin{enumerate}
  \def\labelenumi{\arabic{enumi})}
  \setcounter{enumi}{2}
  \tightlist
  \item
    Mapa de solos e acessibilidade rural do Alto Paraguay (MOPC).
  \end{enumerate}
\item
  \begin{enumerate}
  \def\labelenumi{\arabic{enumi})}
  \setcounter{enumi}{3}
  \tightlist
  \item
    Registro de segurança pública e governança capital regional.
  \end{enumerate}
\end{itemize}}{dist:17_Alto_Paraguay:Fuerte_Olimpo}
\begin{table}[ht!]
\small \begin{tabular}{p{3.0cm}p{0.8cm}p{6.0cm}} \toprule
\textbf{Critério} & \textbf{Nota} & \textbf{Justificativa} \\ \midrule
A - Ameaças Estratégicas & 8.5 & Isolamento tático absoluto; observatório natural em colinas e ausência de alvos militares globais \\
B - Risco Social & 8.5 & Densidade populacional próxima de zero; segurança contra riscos de massa absoluta \\
C - Riscos Naturais & 7.0 & Geologia muito estável e baixo risco de grandes desastres naturais graves \\
D - Autossuficiência & 8.5 & Excepcionais recursos de proteína animal e água fluvial ilimitada \\
E - Institucional & 8.0 & Forte segurança institucional e serviços de saúde de referência capital \\
\midrule \textbf{GSS FINAL} & \textbf{8.1} & \textbf{Seguro (Localidade de elite para refúgio tático extremo, oferecendo o máximo em privacidade e suporte institucional regional).} \\ \bottomrule \end{tabular}
\end{table}
\subsubsection{Vulnerabilidades e
Mitigação}\label{vuln-Fuerte_Olimpo-vulnerabilidades-e-mitigauxe7uxe3o}

\begin{itemize}
\tightlist
\item
  \textbf{Alvo Industrial Estratégico:} A planta da INC é um ponto
  crítico em cenários de sabotagem ou crise nacional (Mitigação:
  residência fora do perímetro imediato da fábrica).
\item
  \textbf{Risco Hídrico:} Inundações sazonais podem afetar a logística
  fluvial e áreas baixas (Mitigação: construção em cotas elevadas,
  monitoramento de avisos do DMH).
\item
  \textbf{Restrições Legais:} Zona de fronteira exige conformidade
  rigorosa para investidores estrangeiros.
\item
  \textbf{Dependência de Serviços:} Necessidade de deslocamento para
  Concepción para serviços especializados.
\end{itemize}

\subsubsection{Dossiê de Campo}\label{dossie-Fuerte_Olimpo-dossiuxea-de-campo}
\begin{description}[style=nextline,leftmargin=0.5cm,font=\bfseries]
\item[Coordenadas]  22°09'S 57°56'W.
\item[Topografia]  Localizada na confluência dos rios Paraguai e Apa. Terreno com elevações calcárias e complexos de cavernas (ex: Santa Caverna). Geologicamente estável.
\item[Alvos Estratégicos]  Industria Nacional del Cemento (INC) - planta vital e estratégica para a construção civil e infraestrutura nacional; Aeródromo de Vallemí (pista asfaltada); Porto de Vallemí (logística fluvial de carga pesada).
\item[Fallout]  Ventos predominantes do Norte/\allowbreak Nordeste.
\item[População]  11.192 (Censo 2022 Final).
\item[Densidade]  \textasciitilde 11,13 hab/\allowbreak km² (Moderada para a região, concentrada no núcleo urbano de Vallemí).
\item[Vias de Acesso]  Ruta PY22 (pavimentação moderna ligando a Concepción). Conexão fluvial estratégica pelos rios Paraguai e Apa. Balsa para o Brasil (Porto Murtinho) nas proximidades.
\item[Serviços]  Hospital da INC e centros de saúde básicos. Dependência de Concepción para alta complexidade.
\item[Custo de Vida]  Médio-baixo; influenciado pela atividade industrial e proximidade com o Brasil.
\item[Preço da Terra]  Lotes urbanos em Vallemí entre US\$ 15.000 e US\$ 40.000; áreas industriais/\allowbreak minerais com valorização elevada.
\item[Hidrologia]  Risco de Inundação Fluvial. Áreas baixas próximas à confluência dos rios Paraguai e Apa são vulneráveis a cheias cíclicas (eventos significativos em 2019 e 2023).
\item[Clima]  Tropical Úmido/\allowbreak Quente. Verões com temperaturas >40°C e tempestades severas.
\item[Energia]  Rede nacional (Itaipu). Subestação dedicada para a INC, conferindo certa prioridade de manutenção, mas sujeita a instabilidades regionais.
\item[Água]  Abundante (Rio Paraguai e Rio Apa).
\item[Qualidade do Solo]  Predomínio de formações calcárias; solo rico em minerais para construção e correção de solos, mas exige manejo para agricultura diversificada.
\item[Recursos Locais]  Maior polo produtor de cimento, cal e mármore do país. Pecuária e pesca comercial.
\item[Segurança]  Cidade industrial com forte presença estatal e corporativa. Zona de fronteira com o Brasil. Presença constante da FTC (Força de Tarefa Conjunta) na região norte do departamento. Taxa de homicídios regional de 14-19/\allowbreak 100k.
\item[Leis Local: Polo de mineração de importância nacional. Restrição de Fronteira (Lei 2532/\allowbreak 05)]  Proibição de compra de terras rurais por estrangeiros de países limítrofes (faixa de 50km).
\end{description}
\textbf{Fonte climática:} NASA POWER Climatology API (período 2001-2020)

\textbf{Fonte luminosa:} estimativa world atlas

\paragraph{Irradiação Solar
(kWh/\allowbreak m²/\allowbreak dia)}\label{irradiauxe7uxe3o-solar-kwhmuxb2dia}

{\def\LTcaptype{none} % do not increment counter
\begin{longtable}[]{@{}
  >{\raggedright\arraybackslash}p{(\linewidth - 24\tabcolsep) * \real{0.0746}}
  >{\raggedright\arraybackslash}p{(\linewidth - 24\tabcolsep) * \real{0.0746}}
  >{\raggedright\arraybackslash}p{(\linewidth - 24\tabcolsep) * \real{0.0746}}
  >{\raggedright\arraybackslash}p{(\linewidth - 24\tabcolsep) * \real{0.0746}}
  >{\raggedright\arraybackslash}p{(\linewidth - 24\tabcolsep) * \real{0.0746}}
  >{\raggedright\arraybackslash}p{(\linewidth - 24\tabcolsep) * \real{0.0746}}
  >{\raggedright\arraybackslash}p{(\linewidth - 24\tabcolsep) * \real{0.0746}}
  >{\raggedright\arraybackslash}p{(\linewidth - 24\tabcolsep) * \real{0.0746}}
  >{\raggedright\arraybackslash}p{(\linewidth - 24\tabcolsep) * \real{0.0746}}
  >{\raggedright\arraybackslash}p{(\linewidth - 24\tabcolsep) * \real{0.0746}}
  >{\raggedright\arraybackslash}p{(\linewidth - 24\tabcolsep) * \real{0.0746}}
  >{\raggedright\arraybackslash}p{(\linewidth - 24\tabcolsep) * \real{0.0746}}
  >{\raggedright\arraybackslash}p{(\linewidth - 24\tabcolsep) * \real{0.1045}}@{}}
\toprule\noalign{}
\begin{minipage}[b]{\linewidth}\raggedright
Jan
\end{minipage} & \begin{minipage}[b]{\linewidth}\raggedright
Fev
\end{minipage} & \begin{minipage}[b]{\linewidth}\raggedright
Mar
\end{minipage} & \begin{minipage}[b]{\linewidth}\raggedright
Abr
\end{minipage} & \begin{minipage}[b]{\linewidth}\raggedright
Mai
\end{minipage} & \begin{minipage}[b]{\linewidth}\raggedright
Jun
\end{minipage} & \begin{minipage}[b]{\linewidth}\raggedright
Jul
\end{minipage} & \begin{minipage}[b]{\linewidth}\raggedright
Ago
\end{minipage} & \begin{minipage}[b]{\linewidth}\raggedright
Set
\end{minipage} & \begin{minipage}[b]{\linewidth}\raggedright
Out
\end{minipage} & \begin{minipage}[b]{\linewidth}\raggedright
Nov
\end{minipage} & \begin{minipage}[b]{\linewidth}\raggedright
Dez
\end{minipage} & \begin{minipage}[b]{\linewidth}\raggedright
Média
\end{minipage} \\
\midrule\noalign{}
\endhead
\bottomrule\noalign{}
\endlastfoot
6.55 & 5.94 & 5.63 & 4.73 & 3.58 & 3.17 & 3.41 & 4.19 & 4.70 & 5.44 &
6.28 & 6.49 & \textbf{5.01} \\
\end{longtable}
}

\textbf{Inclinação solar recomendada:} 22° N (anual) · 32° N (inverno
jun-ago) · 12° N (verão nov-jan)

\paragraph{Precipitação (mm/\allowbreak mês)}\label{precipitauxe7uxe3o-mmmuxeas}

{\def\LTcaptype{none} % do not increment counter
\begin{longtable}[]{@{}
  >{\raggedright\arraybackslash}p{(\linewidth - 24\tabcolsep) * \real{0.0704}}
  >{\raggedright\arraybackslash}p{(\linewidth - 24\tabcolsep) * \real{0.0704}}
  >{\raggedright\arraybackslash}p{(\linewidth - 24\tabcolsep) * \real{0.0704}}
  >{\raggedright\arraybackslash}p{(\linewidth - 24\tabcolsep) * \real{0.0704}}
  >{\raggedright\arraybackslash}p{(\linewidth - 24\tabcolsep) * \real{0.0704}}
  >{\raggedright\arraybackslash}p{(\linewidth - 24\tabcolsep) * \real{0.0704}}
  >{\raggedright\arraybackslash}p{(\linewidth - 24\tabcolsep) * \real{0.0704}}
  >{\raggedright\arraybackslash}p{(\linewidth - 24\tabcolsep) * \real{0.0704}}
  >{\raggedright\arraybackslash}p{(\linewidth - 24\tabcolsep) * \real{0.0704}}
  >{\raggedright\arraybackslash}p{(\linewidth - 24\tabcolsep) * \real{0.0704}}
  >{\raggedright\arraybackslash}p{(\linewidth - 24\tabcolsep) * \real{0.0704}}
  >{\raggedright\arraybackslash}p{(\linewidth - 24\tabcolsep) * \real{0.0704}}
  >{\raggedright\arraybackslash}p{(\linewidth - 24\tabcolsep) * \real{0.1549}}@{}}
\toprule\noalign{}
\begin{minipage}[b]{\linewidth}\raggedright
Jan
\end{minipage} & \begin{minipage}[b]{\linewidth}\raggedright
Fev
\end{minipage} & \begin{minipage}[b]{\linewidth}\raggedright
Mar
\end{minipage} & \begin{minipage}[b]{\linewidth}\raggedright
Abr
\end{minipage} & \begin{minipage}[b]{\linewidth}\raggedright
Mai
\end{minipage} & \begin{minipage}[b]{\linewidth}\raggedright
Jun
\end{minipage} & \begin{minipage}[b]{\linewidth}\raggedright
Jul
\end{minipage} & \begin{minipage}[b]{\linewidth}\raggedright
Ago
\end{minipage} & \begin{minipage}[b]{\linewidth}\raggedright
Set
\end{minipage} & \begin{minipage}[b]{\linewidth}\raggedright
Out
\end{minipage} & \begin{minipage}[b]{\linewidth}\raggedright
Nov
\end{minipage} & \begin{minipage}[b]{\linewidth}\raggedright
Dez
\end{minipage} & \begin{minipage}[b]{\linewidth}\raggedright
Total/\allowbreak ano
\end{minipage} \\
\midrule\noalign{}
\endhead
\bottomrule\noalign{}
\endlastfoot
140 & 144 & 111 & 104 & 101 & 54 & 35 & 20 & 50 & 119 & 192 & 152 &
\textbf{1226 mm} \\
\end{longtable}
}

\paragraph{Poluição Luminosa}\label{poluiuxe7uxe3o-luminosa}

{\def\LTcaptype{none} % do not increment counter
\begin{longtable}[]{@{}ll@{}}
\toprule\noalign{}
Parâmetro & Valor \\
\midrule\noalign{}
\endhead
\bottomrule\noalign{}
\endlastfoot
Escala Bortle & 3 --- Céu rural \\
Radiância artificial & 1.5 nW/\allowbreak cm²/\allowbreak sr \\
\end{longtable}
}
\paragraph{Solo (SoilGrids 2.0, média ponderada 0--30
cm)}

{\def\LTcaptype{none} % do not increment counter
\begin{longtable}[]{@{}ll@{}}
\toprule\noalign{}
Parâmetro & Valor \\
\midrule\noalign{}
\endhead
\bottomrule\noalign{}
\endlastfoot
pH (H$_2$O) & 6.23 \\
Carbono orgânico (SOC) & 13.82 g/\allowbreak kg \\
Argila & 31.72 \% \\
Areia & 36.28 \% \\
Silte (calc.) & 32.0 \% \\
Densidade aparente & 1.36 g/\allowbreak cm³ \\
\textbf{Aptidão agrícola} & \textbf{Alta} \\
\end{longtable}
}

\subsubsection{Indicadores Sociais}

\textbf{Fontes:} PNUD 2020, INE\footnote{Instituto Nacional de Estadística (Instituto Nacional de Estatística), órgão oficial de dados demográficos e censitários.} Censo 2022, INE\footnote{Instituto Nacional de Estadística (Instituto Nacional de Estatística), órgão oficial de dados demográficos e censitários.} EPHC 2023, Ministerio
Público 2024\\
\textbf{Nota:} Valores marcados como \emph{(dept.)} referem-se ao
departamento de Concepción; valores \emph{(dist.)} são específicos deste
distrito. Dados marcados \emph{(est.)} são estimativas.

{\def\LTcaptype{none} % do not increment counter
\begin{longtable}[]{@{}
  >{\raggedright\arraybackslash}p{(\linewidth - 4\tabcolsep) * \real{0.4231}}
  >{\raggedright\arraybackslash}p{(\linewidth - 4\tabcolsep) * \real{0.2692}}
  >{\raggedright\arraybackslash}p{(\linewidth - 4\tabcolsep) * \real{0.3077}}@{}}
\toprule\noalign{}
\begin{minipage}[b]{\linewidth}\raggedright
Indicador
\end{minipage} & \begin{minipage}[b]{\linewidth}\raggedright
Valor
\end{minipage} & \begin{minipage}[b]{\linewidth}\raggedright
Âmbito
\end{minipage} \\
\midrule\noalign{}
\endhead
\bottomrule\noalign{}
\endlastfoot
IDH (2020) & 0.652 & dist. (est.) \\
Ranking IDH nacional & 9/\allowbreak 18 & dept. \\
Esperança de vida & 72,4 anos & dept. \\
Escolaridade média & 8.9 anos & dist. \\
RNB per capita (USD PPA) & 8.400 & dept. \\
Pobreza monetária (\%) & 31,1\% & dept. \\
Pobreza extrema (\%) & 7,8\% & dept. \\
Índice de Gini & 0,460 & dept. \\
Acesso a água potável (\%) & 74,2\% & dept. \\
Acesso a saneamento (\%) & 71,8\% & dept. \\
Taxa de homicídios (est., /\allowbreak 100k hab) & \textasciitilde12--15
(estimativa, acima da média nacional) & dept. \\
Índice de segurança & baixo & dept. \\
População (Censo 2022) & 11.192 hab. & dist. \\
Idade mediana & 26.0 anos & dist. \\
Taxa de fecundidade & 2 & dist. \\
\end{longtable}
}
\subsubsection{Combustível}

\textbf{Referência:} PETROPAR /\allowbreak  postos locais (2024) \textbf{Tipo de
localidade:} Interior

{\def\LTcaptype{none} % do not increment counter
\begin{longtable}[]{@{}lll@{}}
\toprule\noalign{}
Combustível & USD/\allowbreak litro & Gs/\allowbreak litro (aprox.) \\
\midrule\noalign{}
\endhead
\bottomrule\noalign{}
\endlastfoot
Gasolina 93 oct & 1.01 & 7,474 \\
Gasolina 97 oct (premium) & 1.11 & 8,214 \\
Diesel & 0.94 & 6,956 \\
\end{longtable}
}

\begin{quote}
Preços podem variar ±5\% conforme posto e sazonalidade. Chaco e interior
remoto apresentam maior variação.
\end{quote}

\subsubsection{Cobertura Celular}

\textbf{Fonte:} CONATEL PY /\allowbreak  operadoras (2024)

{\def\LTcaptype{none} % do not increment counter
\begin{longtable}[]{@{}ll@{}}
\toprule\noalign{}
Parâmetro & Valor \\
\midrule\noalign{}
\endhead
\bottomrule\noalign{}
\endlastfoot
Cobertura 4G população (dept.) & 78\% \\
Cobertura 4G área rural & 55\% \\
Melhor operadora & Tigo \\
Qualidade rural & moderada \\
\end{longtable}
}

\begin{quote}
Para áreas rurais fora do núcleo urbano, recomenda-se chip Tigo como
principal e Personal como backup.
\end{quote}

\subsubsection{Internet}

\textbf{Fonte:} CONATEL /\allowbreak  Speedtest Ookla (2024)

{\def\LTcaptype{none} % do not increment counter
\begin{longtable}[]{@{}ll@{}}
\toprule\noalign{}
Parâmetro & Valor \\
\midrule\noalign{}
\endhead
\bottomrule\noalign{}
\endlastfoot
Velocidade média download & 38 Mbps \\
Domicílios com internet (dept.) & 48\% \\
Tecnologia predominante & rádio \\
Opção rural & Starlink disponível (\textasciitilde USD 44/\allowbreak mês) \\
\end{longtable}
}

\subsubsection{Mercado Imobiliário e Terra
Rural}

\textbf{Fonte:} INDERT /\allowbreak  Clasificados.com.py (2024)

{\def\LTcaptype{none} % do not increment counter
\begin{longtable}[]{@{}ll@{}}
\toprule\noalign{}
Tipo & Referência \\
\midrule\noalign{}
\endhead
\bottomrule\noalign{}
\endlastfoot
Terra agrícola alta prod. (USD/\allowbreak ha) & 3,500 \\
Imóvel urbano (USD/\allowbreak m²) & 650 \\
Aluguel 2 quartos (USD/\allowbreak mês) & 260 \\
\end{longtable}
}

\begin{quote}
Valores de referência departamental. Localidades menores podem ter
preços 20--40\% abaixo da capital departamental.
\end{quote}

\subsubsection{Saúde}

\textbf{Fonte:} MSPBS /\allowbreak  IPS Paraguay (2024-2026), consolidação
departamental e proxy local

{\def\LTcaptype{none} % do not increment counter
\begin{longtable}[]{@{}ll@{}}
\toprule\noalign{}
Serviço & Disponibilidade \\
\midrule\noalign{}
\endhead
\bottomrule\noalign{}
\endlastfoot
USF /\allowbreak  Posto de Saúde & sim \\
Hospital Regional & sim \\
IPS (seguro social) & sim \\
Farmácia & sim \\
Distância ao hospital de referência & \textasciitilde192 km
(Concepción) \\
\end{longtable}
}

\textbf{Principais estabelecimentos:} USF local; referencia hospitalar
em Concepción

\textbf{Observação para imigrantes:} Atendimento primario local ou em
raio curto; casos de maior complexidade seguem para Concepción.
Cobertura privada continua recomendada para especialidades e urgências
de maior complexidade.
\newpage
\secaoDiagnostico{Puerto Casado}{\mapaDistrito{1702}}{}{dist:17_Alto_Paraguay:Puerto_Casado}
\begin{table}[ht!]
\small \begin{tabular}{p{3.0cm}p{0.8cm}p{6.0cm}} \toprule
\textbf{Critério} & \textbf{Nota} & \textbf{Justificativa} \\ \midrule
A - Ameaças Estratégicas & 7.5 & Isolamento tático estratégico; porto de águas profundas e controle do aqueduto; baixa visibilidade militar \\
B - Risco Social & 8.0 & Densidade populacional extremamente baixa; isolamento social de alto nível \\
C - Riscos Naturais & 6.5 & Geologia estável; nota penalizada pela vulnerabilidade a inundações ribeirinhas \\
D - Autossuficiência & 8.0 & Excelentes recursos de proteína animal e água fluvial estratégica \\
E - Institucional & 7.0 & Presença institucional naval e suporte hídrico tático funcional \\
\midrule \textbf{GSS FINAL} & \textbf{7.4} & \textbf{Seguro (Ideal para quem busca isolamento tático fluvial com abundância de recursos básicos e segurança institucional regional).} \\ \bottomrule \end{tabular}
\end{table}
\subsubsection{Vulnerabilidades e
Mitigação}\label{vuln-Puerto_Casado-vulnerabilidades-e-mitigauxe7uxe3o}

\begin{itemize}
\tightlist
\item
  \textbf{Insegurança Regional:} Atividade insurgente e criminosa na
  fronteira seca (Mitigação: sistemas de segurança redundantes,
  comunicação via satélite e apoio à FTC).
\item
  \textbf{Isolamento Logístico:} Estradas de terra intransitáveis
  (Mitigação: estoque de suprimentos para 60 dias e posse de veículos de
  tração 4x4 pesada).
\item
  \textbf{Instabilidade Energética:} Rede elétrica pouco confiável
  (Mitigação: instalação de sistemas solares off-grid com baterias de
  lítio).
\item
  \textbf{Restrições de Propriedade:} Zona de fronteira exige
  estratégias jurídicas para investidores estrangeiros.
\end{itemize}

\subsubsection{Dossiê de Campo}\label{dossie-Puerto_Casado-dossiuxea-de-campo}
\begin{description}[style=nextline,leftmargin=0.5cm,font=\bfseries]
\item[Coordenadas]  22°31'S 56°52'W.
\item[Topografia]  Relevo ondulado com vastas áreas de pastagens e florestas. Proximidade imediata com o Parque Nacional Paso Bravo (grande reserva de biodiversidade e isolamento). Geologicamente estável.
\item[Alvos Estratégicos]  Grandes estâncias de pecuária de exportação, postos de fronteira seca com o Brasil (Bella Vista).
\item[Fallout]  Ventos predominantes do Norte/\allowbreak Nordeste.
\item[População]  6.153 (Censo 2022 Final).
\item[Densidade]  \textasciitilde 2,74 hab/\allowbreak km² (Baixa densidade, ideal para evitar grandes massas populacionais e garantir privacidade estratégica).
\item[Vias de Acesso]  Acesso Terrestre Crítico. Estradas de terra que sofrem severamente em períodos de chuva. Conexão principal com Concepción e Bella Vista Norte. Isolamento físico frequente.
\item[Serviços]  Infraestrutura básica de saúde (Puestos de Salud) e educação. Dependência crítica de Concepción para alta complexidade.
\item[Custo de Vida]  Reduzido para produtos locais, mas elevado para insumos industrializados devido ao frete.
\item[Preço da Terra]  US\$ 2.000-5.000/\allowbreak ha (áreas de fronteira pecuária e silvicultura).
\item[Hidrologia]  Risco de Isolamento por Chuvas. Transbordamento de arroios tributários do Rio Apa e estradas intransitáveis em períodos de chuva intensa são os principais riscos operacionais.
\item[Clima]  Tropical Úmido. Verões quentes e úmidos com tempestades severas.
\item[Energia]  Rede nacional (ANDE\footnote{Administración Nacional de Electricidade (Administração Nacional de Eletricidade), companhia estatal de energia.}) instável por ser área remota e final de linha; interrupções frequentes.
\item[Água]  Poços artesianos e arroios locais perenes.
\item[Qualidade do Solo]  Solo com boa aptidão para pastagem e agricultura de subsistência; proximidade com áreas de reserva garante biodiversidade.
\item[Recursos Locais]  Pecuária de corte extensiva, exploração florestal sustentável e agricultura familiar. Alto potencial para autossuficiência básica.
\end{description}
\textbf{Fonte climática:} NASA POWER Climatology API (período 2001-2020)

\textbf{Fonte luminosa:} estimativa world atlas

\paragraph{Irradiação Solar
(kWh/\allowbreak m²/\allowbreak dia)}\label{irradiauxe7uxe3o-solar-kwhmuxb2dia}

{\def\LTcaptype{none} % do not increment counter
\begin{longtable}[]{@{}
  >{\raggedright\arraybackslash}p{(\linewidth - 24\tabcolsep) * \real{0.0746}}
  >{\raggedright\arraybackslash}p{(\linewidth - 24\tabcolsep) * \real{0.0746}}
  >{\raggedright\arraybackslash}p{(\linewidth - 24\tabcolsep) * \real{0.0746}}
  >{\raggedright\arraybackslash}p{(\linewidth - 24\tabcolsep) * \real{0.0746}}
  >{\raggedright\arraybackslash}p{(\linewidth - 24\tabcolsep) * \real{0.0746}}
  >{\raggedright\arraybackslash}p{(\linewidth - 24\tabcolsep) * \real{0.0746}}
  >{\raggedright\arraybackslash}p{(\linewidth - 24\tabcolsep) * \real{0.0746}}
  >{\raggedright\arraybackslash}p{(\linewidth - 24\tabcolsep) * \real{0.0746}}
  >{\raggedright\arraybackslash}p{(\linewidth - 24\tabcolsep) * \real{0.0746}}
  >{\raggedright\arraybackslash}p{(\linewidth - 24\tabcolsep) * \real{0.0746}}
  >{\raggedright\arraybackslash}p{(\linewidth - 24\tabcolsep) * \real{0.0746}}
  >{\raggedright\arraybackslash}p{(\linewidth - 24\tabcolsep) * \real{0.0746}}
  >{\raggedright\arraybackslash}p{(\linewidth - 24\tabcolsep) * \real{0.1045}}@{}}
\toprule\noalign{}
\begin{minipage}[b]{\linewidth}\raggedright
Jan
\end{minipage} & \begin{minipage}[b]{\linewidth}\raggedright
Fev
\end{minipage} & \begin{minipage}[b]{\linewidth}\raggedright
Mar
\end{minipage} & \begin{minipage}[b]{\linewidth}\raggedright
Abr
\end{minipage} & \begin{minipage}[b]{\linewidth}\raggedright
Mai
\end{minipage} & \begin{minipage}[b]{\linewidth}\raggedright
Jun
\end{minipage} & \begin{minipage}[b]{\linewidth}\raggedright
Jul
\end{minipage} & \begin{minipage}[b]{\linewidth}\raggedright
Ago
\end{minipage} & \begin{minipage}[b]{\linewidth}\raggedright
Set
\end{minipage} & \begin{minipage}[b]{\linewidth}\raggedright
Out
\end{minipage} & \begin{minipage}[b]{\linewidth}\raggedright
Nov
\end{minipage} & \begin{minipage}[b]{\linewidth}\raggedright
Dez
\end{minipage} & \begin{minipage}[b]{\linewidth}\raggedright
Média
\end{minipage} \\
\midrule\noalign{}
\endhead
\bottomrule\noalign{}
\endlastfoot
6.14 & 5.62 & 5.08 & 4.19 & 3.33 & 2.97 & 3.57 & 4.51 & 4.88 & 5.46 &
6.00 & 5.97 & \textbf{4.81} \\
\end{longtable}
}

\textbf{Inclinação solar recomendada:} 21° N (anual) · 31° N (inverno
jun-ago) · 11° N (verão nov-jan)

\paragraph{Precipitação (mm/\allowbreak mês)}\label{precipitauxe7uxe3o-mmmuxeas}

{\def\LTcaptype{none} % do not increment counter
\begin{longtable}[]{@{}
  >{\raggedright\arraybackslash}p{(\linewidth - 24\tabcolsep) * \real{0.0704}}
  >{\raggedright\arraybackslash}p{(\linewidth - 24\tabcolsep) * \real{0.0704}}
  >{\raggedright\arraybackslash}p{(\linewidth - 24\tabcolsep) * \real{0.0704}}
  >{\raggedright\arraybackslash}p{(\linewidth - 24\tabcolsep) * \real{0.0704}}
  >{\raggedright\arraybackslash}p{(\linewidth - 24\tabcolsep) * \real{0.0704}}
  >{\raggedright\arraybackslash}p{(\linewidth - 24\tabcolsep) * \real{0.0704}}
  >{\raggedright\arraybackslash}p{(\linewidth - 24\tabcolsep) * \real{0.0704}}
  >{\raggedright\arraybackslash}p{(\linewidth - 24\tabcolsep) * \real{0.0704}}
  >{\raggedright\arraybackslash}p{(\linewidth - 24\tabcolsep) * \real{0.0704}}
  >{\raggedright\arraybackslash}p{(\linewidth - 24\tabcolsep) * \real{0.0704}}
  >{\raggedright\arraybackslash}p{(\linewidth - 24\tabcolsep) * \real{0.0704}}
  >{\raggedright\arraybackslash}p{(\linewidth - 24\tabcolsep) * \real{0.0704}}
  >{\raggedright\arraybackslash}p{(\linewidth - 24\tabcolsep) * \real{0.1549}}@{}}
\toprule\noalign{}
\begin{minipage}[b]{\linewidth}\raggedright
Jan
\end{minipage} & \begin{minipage}[b]{\linewidth}\raggedright
Fev
\end{minipage} & \begin{minipage}[b]{\linewidth}\raggedright
Mar
\end{minipage} & \begin{minipage}[b]{\linewidth}\raggedright
Abr
\end{minipage} & \begin{minipage}[b]{\linewidth}\raggedright
Mai
\end{minipage} & \begin{minipage}[b]{\linewidth}\raggedright
Jun
\end{minipage} & \begin{minipage}[b]{\linewidth}\raggedright
Jul
\end{minipage} & \begin{minipage}[b]{\linewidth}\raggedright
Ago
\end{minipage} & \begin{minipage}[b]{\linewidth}\raggedright
Set
\end{minipage} & \begin{minipage}[b]{\linewidth}\raggedright
Out
\end{minipage} & \begin{minipage}[b]{\linewidth}\raggedright
Nov
\end{minipage} & \begin{minipage}[b]{\linewidth}\raggedright
Dez
\end{minipage} & \begin{minipage}[b]{\linewidth}\raggedright
Total/\allowbreak ano
\end{minipage} \\
\midrule\noalign{}
\endhead
\bottomrule\noalign{}
\endlastfoot
129 & 135 & 102 & 63 & 21 & 11 & 7 & 3 & 11 & 39 & 65 & 126 &
\textbf{717 mm} \\
\end{longtable}
}

\paragraph{Poluição Luminosa}\label{poluiuxe7uxe3o-luminosa}

{\def\LTcaptype{none} % do not increment counter
\begin{longtable}[]{@{}ll@{}}
\toprule\noalign{}
Parâmetro & Valor \\
\midrule\noalign{}
\endhead
\bottomrule\noalign{}
\endlastfoot
Escala Bortle & 3 --- Céu rural \\
Radiância artificial & 1.5 nW/\allowbreak cm²/\allowbreak sr \\
\end{longtable}
}
\paragraph{Solo (SoilGrids 2.0, média ponderada 0--30
cm)}

{\def\LTcaptype{none} % do not increment counter
\begin{longtable}[]{@{}ll@{}}
\toprule\noalign{}
Parâmetro & Valor \\
\midrule\noalign{}
\endhead
\bottomrule\noalign{}
\endlastfoot
pH (H$_2$O) & 5.65 \\
Carbono orgânico (SOC) & 14.67 g/\allowbreak kg \\
Argila & 18.83 \% \\
Areia & 61.33 \% \\
Silte (calc.) & 19.8 \% \\
Densidade aparente & 1.31 g/\allowbreak cm³ \\
\textbf{Aptidão agrícola} & \textbf{Média} \\
\end{longtable}
}

\subsubsection{Indicadores Sociais}

\textbf{Fontes:} PNUD 2020, INE\footnote{Instituto Nacional de Estadística (Instituto Nacional de Estatística), órgão oficial de dados demográficos e censitários.} Censo 2022, INE\footnote{Instituto Nacional de Estadística (Instituto Nacional de Estatística), órgão oficial de dados demográficos e censitários.} EPHC 2023, Ministerio
Público 2024\\
\textbf{Nota:} Valores marcados como \emph{(dept.)} referem-se ao
departamento de Concepción; valores \emph{(dist.)} são específicos deste
distrito. Dados marcados \emph{(est.)} são estimativas.

{\def\LTcaptype{none} % do not increment counter
\begin{longtable}[]{@{}
  >{\raggedright\arraybackslash}p{(\linewidth - 4\tabcolsep) * \real{0.4231}}
  >{\raggedright\arraybackslash}p{(\linewidth - 4\tabcolsep) * \real{0.2692}}
  >{\raggedright\arraybackslash}p{(\linewidth - 4\tabcolsep) * \real{0.3077}}@{}}
\toprule\noalign{}
\begin{minipage}[b]{\linewidth}\raggedright
Indicador
\end{minipage} & \begin{minipage}[b]{\linewidth}\raggedright
Valor
\end{minipage} & \begin{minipage}[b]{\linewidth}\raggedright
Âmbito
\end{minipage} \\
\midrule\noalign{}
\endhead
\bottomrule\noalign{}
\endlastfoot
IDH (2020) & 0.628 & dist. (est.) \\
Ranking IDH nacional & 9/\allowbreak 18 & dept. \\
Esperança de vida & 72,4 anos & dept. \\
Escolaridade média & 7.0 anos & dist. \\
RNB per capita (USD PPA) & 8.400 & dept. \\
Pobreza monetária (\%) & 31,1\% & dept. \\
Pobreza extrema (\%) & 7,8\% & dept. \\
Índice de Gini & 0,460 & dept. \\
Acesso a água potável (\%) & 74,2\% & dept. \\
Acesso a saneamento (\%) & 71,8\% & dept. \\
Taxa de homicídios (est., /\allowbreak 100k hab) & \textasciitilde12--15
(estimativa, acima da média nacional) & dept. \\
Índice de segurança & baixo & dept. \\
População (Censo 2022) & 6.153 hab. & dist. \\
Idade mediana & 24.0 anos & dist. \\
Taxa de fecundidade & 3 & dist. \\
\end{longtable}
}
\subsubsection{Combustível}

\textbf{Referência:} PETROPAR /\allowbreak  postos locais (2024) \textbf{Tipo de
localidade:} Interior

{\def\LTcaptype{none} % do not increment counter
\begin{longtable}[]{@{}lll@{}}
\toprule\noalign{}
Combustível & USD/\allowbreak litro & Gs/\allowbreak litro (aprox.) \\
\midrule\noalign{}
\endhead
\bottomrule\noalign{}
\endlastfoot
Gasolina 93 oct & 1.01 & 7,474 \\
Gasolina 97 oct (premium) & 1.11 & 8,214 \\
Diesel & 0.94 & 6,956 \\
\end{longtable}
}

\begin{quote}
Preços podem variar ±5\% conforme posto e sazonalidade. Chaco e interior
remoto apresentam maior variação.
\end{quote}

\subsubsection{Cobertura Celular}

\textbf{Fonte:} CONATEL PY /\allowbreak  operadoras (2024)

{\def\LTcaptype{none} % do not increment counter
\begin{longtable}[]{@{}ll@{}}
\toprule\noalign{}
Parâmetro & Valor \\
\midrule\noalign{}
\endhead
\bottomrule\noalign{}
\endlastfoot
Cobertura 4G população (dept.) & 78\% \\
Cobertura 4G área rural & 55\% \\
Melhor operadora & Tigo \\
Qualidade rural & moderada \\
\end{longtable}
}

\begin{quote}
Para áreas rurais fora do núcleo urbano, recomenda-se chip Tigo como
principal e Personal como backup.
\end{quote}

\subsubsection{Internet}

\textbf{Fonte:} CONATEL /\allowbreak  Speedtest Ookla (2024)

{\def\LTcaptype{none} % do not increment counter
\begin{longtable}[]{@{}ll@{}}
\toprule\noalign{}
Parâmetro & Valor \\
\midrule\noalign{}
\endhead
\bottomrule\noalign{}
\endlastfoot
Velocidade média download & 38 Mbps \\
Domicílios com internet (dept.) & 48\% \\
Tecnologia predominante & rádio \\
Opção rural & Starlink disponível (\textasciitilde USD 44/\allowbreak mês) \\
\end{longtable}
}

\subsubsection{Mercado Imobiliário e Terra
Rural}

\textbf{Fonte:} INDERT /\allowbreak  Clasificados.com.py (2024)

{\def\LTcaptype{none} % do not increment counter
\begin{longtable}[]{@{}ll@{}}
\toprule\noalign{}
Tipo & Referência \\
\midrule\noalign{}
\endhead
\bottomrule\noalign{}
\endlastfoot
Terra agrícola alta prod. (USD/\allowbreak ha) & 3,500 \\
Imóvel urbano (USD/\allowbreak m²) & 650 \\
Aluguel 2 quartos (USD/\allowbreak mês) & 260 \\
\end{longtable}
}

\begin{quote}
Valores de referência departamental. Localidades menores podem ter
preços 20--40\% abaixo da capital departamental.
\end{quote}

\subsubsection{Saúde}

\textbf{Fonte:} MSPBS /\allowbreak  IPS Paraguay (2024-2026), consolidação
departamental e proxy local

{\def\LTcaptype{none} % do not increment counter
\begin{longtable}[]{@{}ll@{}}
\toprule\noalign{}
Serviço & Disponibilidade \\
\midrule\noalign{}
\endhead
\bottomrule\noalign{}
\endlastfoot
USF /\allowbreak  Posto de Saúde & sim \\
Hospital Regional & não \\
IPS (seguro social) & sim \\
Farmácia & sim \\
Distância ao hospital de referência & \textasciitilde737 km
(Concepción) \\
\end{longtable}
}

\textbf{Principais estabelecimentos:} USF local; referencia hospitalar
em Concepción

\textbf{Observação para imigrantes:} Atendimento primario local ou em
raio curto; casos de maior complexidade seguem para Concepción.
Cobertura privada continua recomendada para especialidades e urgências
de maior complexidade.
